\section{Introduction}
\label{sec:intro}
Charged-current pion production by few-GeV neutrinos interacting 
with nuclei (e.g. carbon, oxygen, and argon) is an important process 
for current and future long baseline neutrino oscillation 
experiments~\cite{Abe:2011ks,Ayres:2004js,LBNE}.  
Recent measurements highlight the important role that the nuclear medium plays
in the production and propagation of hadrons produced in neutrino-nucleus
interactions~\cite{MINERVAnuprl,MINERVAnubarprl,mboone-ccqe,mboone-ccqe-nubar}.
These experiments find cross section distortions which are absent in scattering from free
nucleons and affect both event rates and final state kinematics.  
These effects impact oscillation experiments, such as T2K~\cite{t2k-theta13} 
and MiniBooNE~\cite{AguilarArevalo:2010wv}, that rely on the charged-current quasielastic (CCQE)
interaction on bound neutrons, $\nu_\ell n \rightarrow \ell^{-} p$, 
to reconstruct the neutrino energy.  Although this is a relatively 
well-understood reaction with simple kinematics, the reconstruction
and interpretation of events that appear quasielastic 
are complicated by the presence of the nuclear medium.
For example, if a charged-current interaction produces a single $\pi^{+}$ (\cconepiplus),
e.g., $\nu_\ell N(p) \rightarrow \ell^{-} p \pi^{+}$, and the pion is 
absorbed by the target nucleus in a Final State Interaction (FSI), 
the event will mimic the quasielastic topology.
In such a case, the reconstructed neutrino energy may
be significantly underestimated~\cite{Lalakulich:2012hs} and, in the absence of an accurate FSI model, this will lead to a 
bias in the measured oscillation parameters.  Therefore, both
pion production and the effect of the nuclear environment on that 
production must be accurately determined.  

In addition to being absorbed,
pions may undergo elastic, inelastic, or charge-exchange scattering before exiting the nucleus.  Neutrino experiments 
model these processes with Monte Carlo event generators that use particle cascade algorithms constrained by 
cross section measurements of pion
absorption and scattering on various target nuclei.
This technique assumes that interactions of pions created within a nucleus are identical to those of accelerator beam pions, an
assumption which can be probed by measurements of pion production
in electron- and neutrino-scattering experiments.
The only existing electron-scattering experiment on heavy nuclei~\cite{Qian:2009aa} examined 
the ``color transparency'' of pion production, but was done at 
higher energies than those that are relevant to neutrino oscillation experiments;   
hadronic invariant masses (pion kinetic energies) accessed were 
greater than 2.1~GeV (2~GeV).

The earliest neutrino \cconepiplus measurements used hydrogen or deuterium
targets~\cite{Bell:1978fnal,Radecky:1981fn,Kitagaki:1986ct,Allen:1986,WA25:1990} or reported neutrino-nucleon cross 
sections extracted from nuclear target data by model-dependent corrections~\cite{Block:1964,Budagov:1969,SKAT:1989}.  
These data, particularly the ANL~\cite{Radecky:1981fn} and BNL~\cite{Kitagaki:1986ct} data, are used to constrain the 
neutrino-nucleon pion production models contained in event generators, but these constraints are fairly weak because the ANL and BNL 
measurements differ by up to $\sim$40\% in normalization.  A recent re-analysis of the two experiments prefers the 
ANL measurement~\cite{newphil}.

There are a few measurements of $\nu_\mu$ \cconepiplus on nuclear targets, which provide insight into the nuclear effects 
important to neutrino oscillation experiments.  The K2K~\cite{K2K_piprod} and MiniBooNE~\cite{miniboone_piratio} 
collaborations measured the \cconepiplus to CCQE 
cross section ratio on carbon and mineral oil (CH$_{2}$) targets, respectively.  MiniBooNE 
also reported an absolute cross section measurement of \cconepiplus
on a nuclear target (CH$_{2}$) for $E_{\nu}\sim 1$~GeV~\cite{miniboone_piprod}.  This measurement is primarily 
sensitive to pions with kinetic energies from 20 to 400~MeV produced by $\Delta(1232)$ decays.
The kinetic energy spectrum of charged pions reported by MiniBooNE
does not show the suppression of pions predicted by beam-based models of 
FSI~\cite{gibuu-pi,valencia-pi,Rodrigues:2014jfa}, particularly around 160 MeV 
where the total pion-carbon cross section peaks and pion absorption is greatest.  At present, 
oscillation experiments must account for this discrepancy by assigning large 
systematic errors on the size of pion FSI~\cite{t2k_osc_long}.

The analyses presented here measure flux-integrated differential cross sections in pion 
kinetic energy \Tpi and pion angle with respect to the neutrino direction \thetapi.  The signal is defined to be a 
charged-current $\nu_\mu$ interaction in the \minerva tracking detector (mostly CH).  The \cconepi (\ccpi) measurement signal definition 
requires that exactly one (at least one) charged pion exits the target nucleus.  There is no restriction on neutral pions or other mesons.  The 
\cconepi (\ccpi) signal is also restricted to $1.5\le E_\nu \le \unit[10.0]{GeV}$ and hadronic invariant mass $W<1.4$ ($<1.8$) GeV.  Charged-current 
coherent pion production is included in the signal definitions. 

These are the first such measurements on a nuclear target in the few-GeV energy range that is important for the NOvA~\cite{Ayres:2004js} and DUNE~\cite{LBNE} oscillation 
experiments. The \cconepi
measurement is dominated by the excitation of the $\Delta(1232)$ P$_{33}$ resonance, which facilitates comparison 
to theoretical calculations, neutrino event generators, and the MiniBooNE measurement.  The \ccpi measurement, of which the \cconepi events 
are a subset, is complementary since it samples about six resonances and additional nonresonant processes.

The remainder of this paper is organized as follows.  Section~\ref{sec:expt} describes the \minerva experiment and the NuMI beam line.  
Section~\ref{sec:expt_sim} discusses the simulations used to analyze data.  The event reconstruction, including 
track reconstruction, particle identification, and the hadronic recoil energy measurement, is described in Section~\ref{sec:reco}.  The event selection 
criteria for both analyses are provided in Section~\ref{sec:ev_rec}.  Section~\ref{sec:xs} describes the procedure used to extract cross sections 
from the selected events.  Finally, Section~\ref{sec:results} presents 
and discusses the measured cross sections, and Section~\ref{sec:summary} summarizes this paper.



